% !TEX program = pdflatex
% Corvane design notes — ünïcödé comment
\documentclass[11pt,a4paper]{article}
\usepackage[utf8]{inputenc}
\usepackage{amsmath, amssymb}
\usepackage[colorlinks=true]{hyperref}
\importmodule[smglom/sets]{mod/set}
\RequirePackage{xcolor}
\makeatletter
\def\@version{1.2}
\makeatother

\newcommand{\corvane}{\textsc{Corvane}}
\renewcommand*\abstractname{Summary}
\title{A Native Clone of \corvane\ in Rust}
\author{Wasi Master \and Ünïcödé Author}
\date{\today}

\begin{document}
\maketitle

\begin{abstract}
We port 30 syntax modes token for token. Prices: 50\% off, \$10 each, \& more \# of items \{braces\} and under\_score.
Spacing\,tweaks\;and\!negative\/italic\\
line break.
\end{abstract}

\section{Introduction}\label{sec:intro}
As shown in Section~\ref{sec:method} and Equation~\eqref{eq:euler}, see also \cite{knuth1984,lamport94}.
\subsection*{Inline math}
The identity $e^{i\pi} + 1 = 0$ is famous; so is $a_1^2 + b_{12}^{3.5} \leq c$.
Display math: \[ \sum_{k=0}^{n} \binom{n}{k} = 2^n \]
and parenthesized \( x \in \mathbb{R}, |x| < 1 \) too.
$$\int_0^\infty e^{-x^2}\,dx = \frac{\sqrt{\pi}}{2}$$
Mixed: $f(x) = 3.14 x + .5 - 2. * y ; ' " ` ~ # ? @ !$ done.
Math escapes: $\$ \& \% \# \{ \} \_ \, \; \! \/ ok$.
Unclosed math $ x + y
continues onto the next line $ back to text.

\begin{equation}
  \label{eq:euler}
  e^{i\theta} = \cos\theta + i \sin\theta % comment in env
\end{equation}

\begin{itemize}
	\item First item with \emph{emphasis} and \textbf{bold}
	\item Second item: \verb|code| and \texttt{mono}
	\item[$\star$] Custom bullet
\end{itemize}

\begin{table}[htbp]
  \centering
  \begin{tabular}{|l|c|r|}
    \hline
    Mode & Tokens & Status \\ \hline
    zig & 561 & ok \\
    stex & 42pt & 3.5em \\
  \end{tabular}
  \caption{Results}\label{tab:results}
\end{table}

\section
  {Split command arguments}
\ref
{sec:intro} and \label [opt] {lbl}
Stray closing } bracket and ] too.
Numbers 12pt 3.5cm 100% 2e3 and words-with-dashes under_scores.
Non-ASCII: naïve café 日本語 🎉 \'{e} \"{o} \^{a} \~{n} \`{u}
\begin{thebibliography}{9}
\bibitem{knuth1984} D. Knuth, \emph{The \TeX book}, 1984.
\Bibitem{x} and \RBibitem{y}
\end{thebibliography}

\end{document}
{unclosed group \textit{nested {deep} text
\foo@bar{x}[y]{z}